% META: {"id":"1SPE-DERGLOBAL-CO-021","chapitre":"1SPE-DERIVATION-GLOBAL","type_objet":"corrige","exercice_id":"1SPE-DERGLOBAL-EX-021","status":"approved","origin":"generated","status_history":["generated","audited","accepted","approved"],"release_acceptance":"RELEASE_OWNER_BATCH_ACCEPTANCE","acceptance_closure_digest":"sha256:9b3ccf9a81c5520fbb7e03b7b2d3e7bf2057a02834b6d908bf3be5f49f3b553f","acceptance_receipt_digest":"sha256:4fad86f0282e0fa4e0419cca1ad6379ae7af5a280c04a4a90880f90a1c044242"}
% BEGIN-VERIFY
% from sympy import symbols, diff, solve
% x = symbols('x')
% f = x**2 - 4*x + 3
% fp = diff(f, x)
% assert fp == 2*x - 4
% racine = solve(fp, x)
% assert racine == [2]
% assert fp.subs(x, 1) < 0
% assert fp.subs(x, 3) > 0
% END-VERIFY
\begin{corrige}{1SPE-DERGLOBAL-EX-021}
\textbf{1.} On dérive $f(x) = x^2 - 4x + 3$ :
\[
f'(x) = 2x - 4
\]

\textbf{2.} On résout $f'(x) = 0$ :
\[
2x - 4 = 0 \iff x = 2
\]

\textbf{3.} On étudie le signe de $f'(x) = 2(x - 2)$ :
\begin{itemize}
\item Si $x < 2$ : $x - 2 < 0$, donc $f'(x) < 0$ : $f$ est décroissante.
\item Si $x > 2$ : $x - 2 > 0$, donc $f'(x) > 0$ : $f$ est croissante.
\end{itemize}

Tableau de variations de $f$ :

\[
\begin{array}{|c|ccccc|}
\hline
x & -\infty & & 2 & & +\infty \\
\hline
f'(x) & & - & 0 & + & \\
\hline
f & +\infty & \searrow & -1 & \nearrow & +\infty \\
\hline
\end{array}
\]

$f$ admet un \textbf{minimum} en $x = 2$, avec $f(2) = 4 - 8 + 3 = \boxed{-1}$.
\end{corrige}
